\documentclass[11pt]{article}
\usepackage[letterpaper,margin=0.65in,includehead,includefoot,headheight=15pt]{geometry}
\usepackage{amsmath,amssymb}
\usepackage{enumitem}
\usepackage{fancyhdr}
\usepackage{microtype}
\usepackage[hidelinks]{hyperref}
\setlength{\parindent}{0pt}
\setlength{\parskip}{4pt}
\setlist[enumerate]{leftmargin=*,itemsep=4pt,topsep=3pt}
\pagestyle{fancy}
\fancyhf{}
\fancyhead[L]{\textbf{Induction Core — Long-Proof Version}}
\fancyhead[R]{COMP 2300}
\fancyfoot[C]{\thepage}
\renewcommand{\headrulewidth}{0.4pt}

\begin{document}
\textbf{Name:} \rule{3.2in}{0.4pt}\hfill
\textbf{Attempt date:} \rule{1.55in}{0.4pt}

\textbf{Time: 20 minutes.} Write one connected mathematical argument. State
every claim and hypothesis precisely, justify each algebraic step, and make it
clear exactly where the induction hypothesis is used.

\section*{Long Inductive Proof}

\textbf{Definition (Factorial).}
For every integer $n\ge 0$,
\[
  n!=n(n-1)(n-2)\cdots 2\cdot 1,
  \qquad 0!:=1.
\]

\textbf{Definition (Binomial Coefficient).}
For integers $n\ge 0$ and $0\le k\le n$,
\[
  \binom{n}{k}=\frac{n!}{k!(n-k)!}.
\]

Complete both parts as one coherent proof.

\begin{enumerate}[label=\textbf{(\alph*)}]
  \item Prove Pascal's identity
  \[
    \binom{n}{k}=\binom{n-1}{k-1}+\binom{n-1}{k}
    \qquad (n\ge 1,\ 1\le k\le n-1).
  \]
  You may use either a combinatorial argument or algebraic manipulation from
  the factorial definition.

  \item Use Pascal's identity to prove by mathematical induction on $n$ that
  \[
    \sum_{k=0}^{n}\binom{n}{k}=2^n
    \qquad\text{for every }n\ge 0.
  \]
  Your proof must include the proposition $P(n)$, the base case, the induction
  hypothesis, the complete induction step with reindexing, and the conclusion.
\end{enumerate}

\vfill
\textbf{Work:}
\vfill
\vfill
\vfill

\newpage
\textbf{Continuation of work:}
\vfill
\vfill
\vfill
\vfill

\end{document}
